\section{Assignment and Name Binding}

Assignment is the operation that connects source-code names to runtime objects. In Python, this connection is not the same as declaring a typed storage box, as a C programmer might expect. A name is not an integer slot, a floating-point slot, or a character buffer. A name is a binding that allows the program to reach an object.

This section develops that idea through ordinary assignment, reassignment, object sharing, unpacking, augmented assignment, and deletion. The goal is not only to learn the syntax of \texttt{=}, \texttt{+=}, or \texttt{del}, but to understand what happens to names and objects when those statements execute.

The central rule is simple:

\begin{quote}
Python assignment does not copy values into variables. It binds names to objects.
\end{quote}

This rule explains why the same name can later refer to an object of another type, why two names can refer to the same object, why rebinding one name does not automatically affect another name, and why deleting a name is not the same thing as directly destroying an object.

The examples in this section are deliberately small. Assignment is not a decorative part of Python syntax. It is one of the main places where Python's runtime model becomes visible in ordinary code.